\chapter*{报告说明与结论先行}
\addcontentsline{toc}{chapter}{报告说明与结论先行}

\section*{报告边界}

本报告面向研发人员、项目负责人和后续接手者，目标是回答三个实际问题：TwinGuide
当前能做什么、每一步为什么这样设计、继续开发时最容易在哪里出问题。报告以当前
工作区中的源码、病例配置、项目文档、测试和已有输出为唯一资料来源，不引入项目外
文献，也不把临床常识扩写为项目已经实现的功能。

所谓“完整阅读”，在本报告中指：盘点并阅读 \projectpath{src/twin_guide} 下全部
Python 模块，核对 \projectpath{docs}、\projectpath{tests}、\projectpath{examples}
与 \projectpath{data/cases}，并检查 tooth-47 的 JSON/YAML、过程报告、过程图和最终
STL。对 STL/PLY/GLB 等二进制网格采用拓扑、尺寸、体积和可视化检查，而非无意义的
逐字节解释。

\section*{结论先行}

TwinGuide 不是一个单纯“把几个 STL 做并集”的脚本。它已经形成了由
\term{病例语义、纯几何规划、Blender 实体化和独立结果检查}组成的分层流程：

\begin{enumerate}
  \item 从导套装配体的连通分量中识别真正具有轴孔的两根导管，并确定轴线、上下端和 C 口方向；
  \item 对牙列执行 FDI 牙位识别和导板映射，建立可用于“U 侧、背 U 侧、近远中和牙合方向”的局部语义坐标；
  \item 规划导孔、双导管操作窗和按牙位范围生成的轴扫掠观察窗；
  \item 在真实导管外壁与原始导板外壁上选择连接锚点，生成四根平滑主连接梁和可选 Y 型按压梁；
  \item 通过牙列净距裁剪、体素融合、精确差集、功能孔复切和手机运动包络差集，形成最终封闭 STL；
  \item 重新构建几何计划，对最终 STL 的拓扑、导管保留、梁连续性、孔道和观察窗进行检查。
\end{enumerate}

当前 tooth-47 是缺失下颌右侧第二磨牙 47 的黄金回归病例。现有结果 STL 为单一封闭
体，含 138,861 个顶点和 277,774 个三角面，体积约 $6670.72\,\mathrm{mm^3}$；观察窗
过程报告的全部 QA 项为真，手机包络也为封闭体。

\begin{takeaway}
项目的核心价值是把“牙位语义”和“几何建模”连成了可审计的数据流；当前最大工程风险
则是版本—产物追溯仍不充分、多个阶段仍标记为 experimental、布尔/体素过程成本高，
且手机避让尚未进入独立验证器的完整复算。
\end{takeaway}

